\section{Yadro reaksiyalari}



\begin{enumerate}

\item Quyidagi zarrachalarning qaysi biri yadroga oson kira oladi va yadro reaksiyasini yuzaga chiqaradi?
\begin{enumerate}[label=\Alph*)]
\item neytron.
\item proton.
\item $\alpha$-zarracha.
\item elektron.
\end{enumerate}

\item Proton, antiproton va $\gamma$-kvant berilgan. Elektr maydon ularning qaysilariga tezlanish beradi?
\begin{enumerate}[label=\Alph*)]
\item proton va antiprotonga.
\item antiproton va $\gamma$-kvantga.
\item proton va $\gamma$-kvantga.
\item faqat protonga.
\end{enumerate}

\item Quyidagi yadro reaksiyasidagi noma'lum zarracha nima?
$$^7_3\text{Li} + ^1_1\text{p} \rightarrow ^7_4\text{Be} + x. \quad x - ?$$
\begin{enumerate}[label=\Alph*)]
\item gamma kvant.
\item elektron.
\item neytron.
\item pozitron.
\end{enumerate}

\item Texnetsiyni sun'iy usulda olishda qanday zarra hosil bo'ladi?
$$^{94}_{42}\text{Mo} + ^2_1\text{H} \rightarrow ^{95}_{43}\text{Tc} + x.$$
\begin{enumerate}[label=\Alph*)]
\item elektron.
\item proton.
\item foton.
\item neytron.
\end{enumerate}

\item Quyidagi yadroviy reaksiyadagi noma'lum zarrani aniqlang?
$$^{19}_9\text{F} + ^1_1\text{p} \rightarrow ^{16}_8\text{O} + x.$$
\begin{enumerate}[label=\Alph*)]
\item elektron.
\item $^{17}_8\text{O}$.
\item proton.
\item $^4_2\text{He}$.
\end{enumerate}

\setcounter{enumi}{5}
\item Quyidagi yadro reaksiyada hosil bo'ladigan elementni aniqlang?
$$^{14}_7\text{N} + ^1_0\text{n} \rightarrow x + ^4_2\text{He}$$
\begin{enumerate}[label=\Alph*)]
\item $^4_2\text{He}$.
\item $^{10}_5\text{B}$.
\item $^{16}_8\text{O}$.
\item $^{11}_5\text{B}$.
\end{enumerate}

\item Quyidagi yadro reaksiyasidagi noma'lum yadroni aniqlang:
$$^{14}_7\text{N} + ^1_0\text{n} \rightarrow x + ^1_1\text{p}$$
\begin{enumerate}[label=\Alph*)]
\item $^{14}_6\text{C}$.
\item $^{16}_7\text{N}$.
\item $^{15}_7\text{N}$.
\item $^{15}_6\text{C}$.
\end{enumerate}

\item $^{27}_{13}\text{Al} + ^1_0\text{n} \rightarrow ^4_2\text{He} + x$ reaksiya natijasida qanday element hosil bo'ladi?
\begin{enumerate}[label=\Alph*)]
\item $^{24}_{12}\text{Mg}$.
\item $^{23}_{12}\text{Mg}$.
\item $^{23}_{11}\text{Na}$.
\item $^{24}_{11}\text{Na}$.
\end{enumerate}

\item Ushbu reaksiya natijasida qanday element hosil bo'ladi?
$$^{54}_{26}\text{Fe} + ^1_0\text{n} \rightarrow ^1_1\text{H} + x$$
\begin{enumerate}[label=\Alph*)]
\item $^{54}_{25}\text{Mn}$.
\item $^{56}_{26}\text{Fe}$.
\item $^{55}_{25}\text{Mn}$.
\item $^{52}_{24}\text{Cr}$.
\end{enumerate}

\item Quyidagi og'ir yadroning bo'linish reaksiyasida qanday zarracha ajralib chiqadi?
$$^{235}_{92}\text{U} + ^1_0\text{n} \rightarrow ^{135}_{52}\text{Te} + ^{99}_{40}\text{Zr} + x$$
\begin{enumerate}[label=\Alph*)]
\item 2 ta proton.
\item 1 ta proton va 1 ta neytron.
\item 2 ta neytron.
\item 3 ta neytron.
\end{enumerate}

\item Quyidagi yadroviy reaksiyadagi noma'lum zarrani aniqlang.
$$^{198}_{80}\text{Hg} + ^1_0\text{n} \rightarrow ^{198}_{79}\text{Au} + x$$
\begin{enumerate}[label=\Alph*)]
\item proton.
\item elektron.
\item neytron.
\item $\alpha$-zarra.
\end{enumerate}

\item $^{80}_{35}\text{Br}$ brom izotopini neytron bilan parchalaganda $\alpha$-zarra va yana bir element hosil bo'ladi. U qanday element?
\begin{enumerate}[label=\Alph*)]
\item $^{77}_{33}\text{As}$.
\item $^{73}_{32}\text{Ge}$.
\item $^{79}_{34}\text{Se}$.
\item $^{70}_{31}\text{Ga}$.
\end{enumerate}

\item $^{24}_{12}\text{Mg}$ magniy yadrosi neytron yutib, $^{24}_{11}\text{Na}$ ga aylanadi. Bu reaksiyada yana qanday zarra hosil bo'ladi?
\begin{enumerate}[label=\Alph*)]
\item elektron.
\item foton.
\item $\alpha$-zarra.
\item proton.
\end{enumerate}

\setcounter{enumi}{13}
\item Quyidagi reaksiyada qanday zarracha ajralib chiqadi?
$$^{14}_7\text{N} + ^4_2\text{He} \rightarrow ^{17}_8\text{O} + x$$
\begin{enumerate}[label=\Alph*)]
\item elektron.
\item pozitron.
\item proton.
\item neytron.
\end{enumerate}

\item Quyidagi reaksiya natijasida qanday zarra hosil bo'ladi?
$$^{27}_{13}\text{Al} + ^4_2\text{He} \rightarrow ^{30}_{15}\text{P} + x$$
\begin{enumerate}[label=\Alph*)]
\item $\alpha$-zarra.
\item elektron.
\item pozitron.
\item neytron.
\end{enumerate}

\item $^{27}_{13}\text{Al} + ^4_2\text{He} \rightarrow x + ^1_0\text{n}$ reaksiya natijasida neytrondan tashqari yana qaysi yadro hosil bo'ladi?
\begin{enumerate}[label=\Alph*)]
\item $^{31}_{16}\text{S}$.
\item $^{30}_{15}\text{P}$.
\item $^{27}_{14}\text{Si}$.
\item $^{31}_{15}\text{P}$.
\end{enumerate}

\item Quyidagi yadroviy reaksiyadagi noma'lum elementni aniqlang?
$$^{27}_{13}\text{Al} + ^4_2\text{He} \rightarrow x + ^1_1\text{H}$$
\begin{enumerate}[label=\Alph*)]
\item $^{23}_{11}\text{Na}$.
\item $^{14}_{7}\text{N}$.
\item $^{30}_{14}\text{Si}$.
\item $^{31}_{15}\text{P}$.
\end{enumerate}

\item Neon $^{20}_{10}\text{Ne}$ izotopini $\alpha$-zarra bilan bombardimon qilganda, proton va noma'lum element hosil bo'ladi. U qanday element?
\begin{enumerate}[label=\Alph*)]
\item $^{21}_{10}\text{Ne}$.
\item $^{24}_{12}\text{Mg}$.
\item $^{25}_{11}\text{Na}$.
\item $^{23}_{11}\text{Na}$.
\end{enumerate}

\item Quyidagi yadro reaksiyasida qanday yangi yadro hosil bo'ladi?
$$^{253}_{99}\text{Es} + ^4_2\text{He} \rightarrow x + ^1_0\text{n}$$
\begin{enumerate}[label=\Alph*)]
\item $^{256}_{100}\text{Fm}$.
\item $^{260}_{104}\text{Ku}$.
\item $^{256}_{101}\text{Md}$.
\item $^{247}_{96}\text{Cm}$.
\end{enumerate}

\item $^{9}_4\text{Be} + ^2_1\text{H} \rightarrow ^{10}_5\text{B} + x$ yadro reaksiyasi natijasida ajralib chiquvchi zarrachani toping.
\begin{enumerate}[label=\Alph*)]
\item $\alpha$-zarracha.
\item $^1_0\text{n}$.
\item $e^-$.
\item $\gamma$.
\end{enumerate}

\item Reaksiyada yetishmaydigan zarrachani aniqlang.
$$^{56}_{26}\text{Fe} + ^2_1\text{H} \rightarrow ^{54}_{25}\text{Mn} + x$$
\begin{enumerate}[label=\Alph*)]
\item pozitron.
\item $\alpha$-zarra.
\item elektron.
\item neytron.
\end{enumerate}

\item Quyidagi reaksiyada deyteriy nima bilan bombardimon qilinganligi aniqlansin:
$$^2_1\text{H} + x \rightarrow ^1_1\text{H} + ^1_0\text{n}$$
\begin{enumerate}[label=\Alph*)]
\item $\gamma$-kvant.
\item $e$-elektron.
\item $n$-neytron.
\item $\alpha$-zarra.
\end{enumerate}

\setcounter{enumi}{22}
\item Quyidagi yadroviy reaksiyada qanday zarra ajraladi?
$$^{13}_6\text{C} + ^0_1\text{e} \rightarrow ^{12}_6\text{C} + x$$
\begin{enumerate}[label=\Alph*)]
\item antiproton.
\item $\alpha$-zarra.
\item neytron.
\item proton.
\end{enumerate}

\item Uran izotopining bitta yadrosi bo'linganda, $200$ MeV energiya ajraladi. Shu izotop $470$ grammining barcha yadrolari bo'linganda, qanday (kJ) energiya ajraladi? ($1\text{ MeV}=1,6\cdot 10^{-13}$ J).
\begin{enumerate}[label=\Alph*)]
\item $3,85\cdot 10^{10}$.
\item $3,85\cdot 10^{11}$.
\item $9,5\cdot 10^{11}$.
\item $4,7\cdot 10^{11}$.
\end{enumerate}

\item Ushbu moddalarning qaysi biri yadro reaktorlarida issiqlik tashuvchi modda sifatida ishlatiladi?
1) uran; 2) grafit; 3) kadmiy; 4) oddiy suv; 5) plutoniy; 6) suyuq natriy.
\begin{enumerate}[label=\Alph*)]
\item 1.
\item 2.
\item 3 va 4.
\item 4 va 6.
\end{enumerate}

\item Quyida sanalgan moddalardan qaysilari yadro reaktorlarida neytronlarni yutuvchi moddalar sifatida ishlatiladi?
1) uran; 2) grafit; 3) kadmiy; 4) og'ir suv; 5) bor; 6) plutoniy.
\begin{enumerate}[label=\Alph*)]
\item 2, 4.
\item 3, 5.
\item 2, 3.
\item 3, 4.
\end{enumerate}

\item Quyida sanalgan moddalardan qaysilari yadro reaktorlarida yadro yoqilg'isi sifatida ishlatiladi?
1) uran; 2) grafit; 3) kadmiy; 4) og'ir suv; 5) bor; 6) plutoniy.
\begin{enumerate}[label=\Alph*)]
\item 1 va 5.
\item 1 va 2.
\item 1 va 3.
\item 1 va 6.
\end{enumerate}

\item Neytronlarning ko'payish koeffitsiyenti $k=1,7$ ga teng bo'lsa reaktor qanday ishlaydi?
\begin{enumerate}[label=\Alph*)]
\item normal rejimda kechadi.
\item reaksiya to'xtaydi.
\item reaksiya normadan tezlashib ketib, portlash yuzaga keladi.
\item reaksiyaning borishi $k$ ga bog'liq emas.
\end{enumerate}

\setcounter{enumi}{28}
\item Ushbu $^2_1\text{H} + ^3_1\text{H} \rightarrow ^4_2\text{He} + ^1_0\text{n}$ termoyadro reaksiyasida $17,4$ MeV energiya ajraladi. $4$ g geliy sintezida necha MeV energiya ajraladi? ($N_A=6\cdot 10^{23}\text{ mol}^{-1}$, $M=4\cdot 10^{-3}\text{ kg/mol}$).
\begin{enumerate}[label=\Alph*)]
\item $1,1\cdot 10^{24}$.
\item $11\cdot 10^{25}$.
\item $1,8\cdot 10^{25}$.
\item $1\cdot 10^{25}$.
\end{enumerate}

\item Quyidagi yadroviy reaksiyada ajraladigan energiyani aniqlang (MeV):
$$^6_3\text{Li} + ^2_1\text{H} \rightarrow ^4_2\text{He} + ^4_2\text{He}$$
$\text{M}(^6_3\text{Li})=6,01513\text{ u}$, $\text{M}(^2_1\text{H})=2,01410\text{ u}$, $\text{M}(^4_2\text{He})=4,0026\text{ u}$, $1\text{ u}=1,66\cdot 10^{-27}\text{ kg}$.
\begin{enumerate}[label=\Alph*)]
\item 22,4.
\item 40,1.
\item 50.
\item 66,8.
\end{enumerate}

\item $^{19}_9\text{F}$ va $^1_1\text{p}$ massalarining yig'indisi, $^{16}_8\text{O}$ va $^4_2\text{He}$ massalarining yig'indisidan katta. Quyidagi yadro reaksiyasida energiya ajraladimi yoki yutiladimi?
$$^{19}_9\text{F} + ^1_1\text{p} \rightarrow ^{16}_8\text{O} + ^4_2\text{He}$$
\begin{enumerate}[label=\Alph*)]
\item bunday reaksiya amalga oshmaydi.
\item ajraladi.
\item energiya ajralmaydi va yutilmaydi.
\item yutiladi.
\end{enumerate}

\item Quyidagi termoyadro reaksiyasida necha MeV energiya ajraladi?
$$^2_1\text{H} + ^2_1\text{H} \rightarrow ^4_2\text{He} + \Delta\text{E}$$
$^2_1\text{H}$ deytronning bog'lanish energiyasi $2,2$ MeV.
$^4_2\text{He}$ geliyning bog'lanish energiyasi $28,3$ MeV.
\begin{enumerate}[label=\Alph*)]
\item 30,5.
\item 26,1.
\item 20.
\item 23,9.
\end{enumerate}

\item Quyida keltirilganlarning qaysi biri annigilatsiya hodisasini aks ettiradi?
\begin{enumerate}[label=\Alph*)]
\item $^{27}_{13}\text{Al} + \gamma = ^{26}_{12}\text{Mg} + ^1_1\text{H}$.
\item $e + \tilde{e} = \gamma$.
\item $e + \tilde{e} = 2\gamma$.
\item $n \rightarrow ^1_1\text{H} + e + \nu$.
\end{enumerate}

\end{enumerate}





\begin{table}[!ht]
\centering
\begin{tabular}{|c|c|c|c|c|c|c|c|c|c|c|}
\hline
1a & 2a & 3c & 4d & 5d & 6d & 7a & 8d & 9a & 10c & 11a \\ \hline
12a & 13d & 14c & 15d & 16b & 17c & 18d & 19c & 20b & 21b & 22a \\ \hline
23d & 24a & 25d & 26b & 27d & 28c & 29d & 30a & 31b & 32d & 33c \\ \hline
\end{tabular}
\end{table}